\section{Dataset}
\label{sec:dataset}

Buscamos un dataset público que tuviera voces tanto reales como clonadas y que ademas fueran por pares (i.e. un audio de voz real cuya transcripción es "Hoy he comido en un restaurante chino" y un audio clonando dicha voz diciendo lo mismo). Esto debido a que planteamos hacer explicabilidad con contraejemplos con los pares, pero al final no pudimos. Además, priorizamos la búsqueda de un dataset multilingüe para comparar las diferentes características en diversos idiomas, así como las diferencias en predicción y explicación.

Vamos a comentar los que hemos tenido en cuenta:

\begin{enumerate}
	\item HABLA~\parencite{TamayoFlorez2023}: Este dataset cuenta con una gran cantidad de audios tanto reales, como generados y voces clonadas.
	\item audioMNIST~\parencite{Becker2024}: Este dataset cuenta con 30 mil audios reales de los números del 0 al 9 en inglés.
	\item ASVSpoof2021~\parencite{asvspoof2021_zenodo}: Este dataset forma parte de una serie de desafios internacionales organizados por la comunidad de verificación automática de audio cuyo objetivo es evaluar y mejorar los sistemas capaces de detectar voz manipulada. El dataset se compone de tres grupos de audio, todos ellos con voz real y generada mezclada, sin pares. 
	\item Mozilla Common Voice~\parencite{MozillaCommonVoice}: Este dataset cobntiene voces de personas reales en varios idiomas, aunque sin voces clonadas.
	\item Wavefake~\parencite{wavefakeKaggle}: Este dataset de Kaggle cuenta con gran cantidad de audios generados, pero sin audios reales.
	\item In-the-wild:~\parencite{2022In-the-Wild}: Este dataset busca condiciones realistas recopilando audios reales y clonados en inglés.
\end{enumerate}

Finalmente, el dataset por el que nos decantamos fue ODSS~\parencite{2023ODSS}. Hemos usado este conjunto de audios porque ya tiene creado un par para cada audio, lo que nos ahorra tiempo de generar nuestros propios audios. A esta ventaja se le suma que ODSS ya contenía audios en inglés y en español, ambos idiomas son sumamente hablados por lo que es interesante centrarse en ellos.

De ODSS hemos tomado un subconjunto de audios (200 en total) entre los que tenemos 50 pares en español y 50 en inglés, por lo que 100 en cada idioma.

\subsection*{Análisis del dataset}

De este subconjunto de audios extraídos sacaremos todas sus características que explicaremos más adelante junto con su espectrograma, su duración, su categoría (natural o generado) y el audio. Esto hace un total de 47 columnas por cada audio ya que de características como las formantes y los MFCCs que no solo se saca su valor promedio sino que también se sacará su distribución a lo largo del audio. Un resumen de los tipos de datos que encontramos en nuestro dataset sería:

\begin{itemize}
  \item \textbf{Datos tabulares}: La mayoría de los datos sobre las características del audio.
  \item \textbf{Series temporales}: Muestran la evolución de las formantes o los MFCCs a lo largo del audio.
  \item \textbf{Datos de audio}: Se corresponde con el archivo de audio (en la práctica es la ruta al archivo de audio).
  \item \textbf{Datos de texto}: Se corresponde con las transcripciones de los audios.
\end{itemize}